
% Author       : Yulong Wang
% Date         : 2026-08-30 10:21:21
% LastEditors  : Yulong Wang
% LastEditTime : 2026-08-30 10:21:22
% FilePath     : /.agents/home/wyl/project/knowledge/political_science/central_local_relations/central_local_relations_report.tex
% Description  :

% !TeX program = xelatex
% ==============================================================================
% 央地关系报告:权力流向与形态的二维分析
% 整理自领域包 political_science/central_local_relations/
% 编译方式: xelatex central_local_relations_report.tex (需 TeX Live / CTeX 宏集)
% ==============================================================================
\documentclass[UTF8,12pt,a4paper]{ctexrep}

\usepackage[a4paper,margin=2.4cm]{geometry}
\usepackage{amsmath,amssymb}
\usepackage{booktabs}
\usepackage{tabularx}
\usepackage{array}
\usepackage{caption}
\usepackage{enumitem}
\usepackage{xcolor}
\usepackage[hidelinks]{hyperref}
\usepackage{url}
\urlstyle{tt}

\captionsetup{font=small,labelfont=bf}
\setlist{nosep,itemsep=0.15em}
\setlength{\parskip}{0.35em}
\hypersetup{
  pdftitle  = {央地关系报告:权力流向与形态的二维分析},
  pdfauthor = {整理自 central\_local\_relations 领域包}
}

% ---- 证据状态标注 ----------------------------------------------------------
% 在正文中用 \src{...} 标注证据状态, 图例见附录 B
\newcommand{\src}[1]{\textsuperscript{[\emph{#1}]}}

\begin{document}

% ==============================================================================
\title{{\bfseries 央地关系报告:权力流向与形态的二维分析}\\[0.6em]
  \large 整理自领域包 \texttt{political\_science/central\_local\_relations/}}
\author{整理:Yulong Wang}
\date{2026年7月}
\maketitle

\begin{abstract}
央地关系(central--local relations)是政治学中最古老也最现代的问题之一。从秦始皇废分封设郡县,到当代中国的条块治理与美国的联邦--州博弈,所有大型政治体都面对同一个结构性张力:如何在保证中央权威的同时,赋予地方有效治理所需的自主权。本报告把 \texttt{central\_local\_relations/} 领域包中分散的笔记、概念辨析、历史分析与当代案例整合为一份连贯的专题报告:第一章说明央地关系为什么重要;第二章提出以``权力流向''与``权力形态''两个维度重新刻画集权与分权的分析框架;第三章至第五章分别用这一框架分析中国古代(唐、宋、明)、美国联邦制与中国当代的央地关系;第六章汇总开放性问题。

本报告遵循包内的证据分层约定:可查证事实、工作定义与制度推理分别标注(图例见附录 B),凡属推测之处均明确声明,不把推测写成确定结论。报告不涉及对当代中国政治的规范性评价。
\end{abstract}

\tableofcontents

% ------------------------------------------------------------------------------
\section*{报告体例说明}
\addcontentsline{toc}{section}{报告体例说明}

{\sloppy
\begin{enumerate}[label=(\arabic*)]
  \item \textbf{材料来源}。本报告整合自领域包内下列文件:种子笔记 \url{conflict.md};总述 \url{01_research_overview.md};概念辨析 \url{11_term_disambiguation.md};领域地图 \url{12_field_map.md};边界罗盘 \url{02_frontier_compass.md};论断--证据台账 \url{17_claims_evidence_ledger.md};研究缺口表 \url{20_gap_table.md};候选项目 \url{22_project_candidates.md};问答缓冲区 \url{90_QA.md};案例研究 \url{case_studies/southern_power_grid.md} 与 \url{case_studies/land_finance.md}。完整清单见附录 A。
  \item \textbf{证据标注}。正文中用 \src{SC}(source\_checked,公开可查证)、\src{WD}(working\_definition,包内工作定义)、\src{IP}(interpretation,基于制度安排的反向推理)、\src{NV}(NEEDS\_VERIFICATION,待核验)标注关键论断的证据状态,图例与汇总见附录 B。
  \item \textbf{引用说明}。文中出现的作者与年份均来自包内已有记录;包内标注``NEEDS\_VERIFICATION''的条目在正文中注明,未逐条重新核验。本报告不新增任何未经包内记录的引文(见附录 C)。
  \item \textbf{术语约定}。关键术语首次出现时附英文;``行政分权''与``行动集权''是本包的工作定义,与标准文献用法的差异在 2.4 节说明。
\end{enumerate}
}

% ==============================================================================
\chapter{央地关系的重要性}

\section{治理的第一性原理:规模、授权与代理}

任何超过一定规模的政治体,都面临一个无法回避的事实:\textbf{统治者不可能亲自处理所有事务}。他必须授权——把一部分决策权和执行权交给代理人(地方官、将领、部门首长)。而一旦授权,就产生了政治学中最经典的问题:代理人可能不按委托人的意愿行事。这就是委托--代理问题(principal--agent problem)在空间治理上的形态:中央是委托人,地方是代理人,核心障碍是信息不对称(information asymmetry)与激励不相容(incentive incompatibility)。

由此可以得到一个反直觉的结论:\textbf{央地矛盾不是某个特定制度或文化的产物,而是``规模治理''本身的结构性副产品}。小国寡民不需要讨论央地关系——因为根本就没有``地方''这个层级。国家越大、层级越多、治理任务越复杂,央地之间的信息损耗、利益分歧与监督成本就越高,央地矛盾也就越难以回避。

\section{历史上的学费}

人类为理解央地关系付出了极其高昂的学费。表~\ref{tab:tuition} 列出了几笔最著名的``学费'':

\begin{table}[htbp]
  \centering
  \small
  \caption{历史上央地关系失衡的代价}
  \label{tab:tuition}
  \begin{tabularx}{\linewidth}{@{}lXX@{}}
    \toprule
    \textbf{案例} & \textbf{央地关系病症} & \textbf{代价} \\
    \midrule
    罗马帝国 & 从共和制到元首制再到四帝共治,本质上是一次又一次央地权力重组实验 & 最终以帝国分裂告终 \\
    唐朝 & 节度使制度下的行动集权失控 & 安史之乱(755--763)将当时世界最强大的帝国从巅峰打落 \\
    宋朝 & 行政分权过度导致军事效能崩溃 & 以岁币换和平、以积弱换稳定,对外屡战屡败 \\
    美国 & 州能否脱离联邦的问题悬而未决 & 南北战争(1861--1865)以约62万人的生命作答 \\
    \bottomrule
  \end{tabularx}
\end{table}

这些案例的共同点是:它们都不是``小失误'',而是足以改变文明走向的制度性失败。央地关系的设计失误,其代价以王朝的兴衰甚至千万人的生命计。

\section{当代相关性}

央地关系不是博物馆里的问题,它每天都在当代政治中重演:

\begin{itemize}
  \item \textbf{2024 年初,美国得克萨斯州与联邦政府在边境移民问题上武装对峙}——联邦制的央地矛盾以近乎冲突的方式显形;
  \item \textbf{COVID-19 大流行中,各国联邦与地方政府的协调或失调}——检验了央地权力配置在危机中的真实效能;
  \item \textbf{中国在经济增长与风险防控之间不断调整条块权力配置}——2021 年以来土地财政收缩引发的地方财政压力,是央地财政关系紧张的当代表现(见 5.3 节)。
\end{itemize}

\section{一个语言问题:为什么需要新的分析框架}

在研究央地关系之前,需要先解决一个语言障碍。日常话语中``集权''与``分权''被当作一根轴上的两个方向:权力流向中央叫集权,流向地方叫分权。但这个一维叙事存在系统性混淆:

\begin{itemize}
  \item 宋朝被称作``中央集权''的典范——但它的实际动作是把执行权拆散到互不统属的机构;
  \item 唐朝节度使被称作``地方分权''——但它的实际动作是把军、财、政权力打包给同一个人。
\end{itemize}

两者被放在同一根轴上讨论,却在结构上南辕北辙。这说明:\textbf{``集权/分权''的话语同时混用了两个不同的维度——权力的流向与权力的形态}。只有把这两个维度分开,才能准确描述唐、宋、明、美国联邦制与中国当代治理各自的真实结构。这就是第二章要建立的二维分析框架。

% ==============================================================================
\chapter{概念分析:关于权力的流向与形态的二维划分}

\section{传统话语的混淆}

先看两个熟悉的例子,它们将贯穿本报告:

\begin{itemize}
  \item \textbf{宋朝}。权力流向中央(``强干弱枝''),但中央``收''权力的方式是把前线的执行权拆散:统兵的三衙、调兵的枢密院、监军的文官互不统属。结果前线指挥员没有任何一个拥有完成任务的完整工具箱。从``流向''看这是集权,从``形态''看却是执行权的分散拆解。
  \item \textbf{唐朝}。权力流向地方(节度使),但``放''权力的方式是把军事指挥、地方财政、行政人事打包给同一个人。从``流向''看这是分权,从``形态''看却是执行权的集中打包。
\end{itemize}

表~\ref{tab:mixup} 总结了这一混淆:

\begin{table}[htbp]
  \centering
  \small
  \caption{传统表述与实际动作的错位}
  \label{tab:mixup}
  \begin{tabularx}{\linewidth}{@{}p{3.2cm}XX@{}}
    \toprule
    \textbf{传统表述} & \textbf{实际动作} & \textbf{权力形态} \\
    \midrule
    ``中央集权''(如宋朝) & 通过相互制衡将最高权力收归中央 & 实质是行政分权:执行权被拆解 \\
    ``地方分权''(如唐朝节度使) & 把军事、财政、行政权力下放地方 & 实质是行动集权:前线权力集中打包 \\
    \bottomrule
  \end{tabularx}
\end{table}

\section{两个维度与二维矩阵}

把``集权/分权''从一维扩展为二维,需要分别问两个问题:

\begin{itemize}
  \item \textbf{维度一:权力流向(决策权归属)}。最终拍板权、人事任免权、财政审批权归谁?流向中央,还是流向地方?
  \item \textbf{维度二:权力形态(执行权分布)}。执行某类任务所需的权力,是分散在不同部门互相制衡,还是集中打包给同一主体?
\end{itemize}

两个维度交叉,得到表~\ref{tab:matrix} 的二维矩阵:

\begin{table}[htbp]
  \centering
  \small
  \caption{权力流向 $\times$ 权力形态的二维矩阵}
  \label{tab:matrix}
  \begin{tabularx}{\linewidth}{@{}Xcc@{}}
    \toprule
    & \multicolumn{2}{c}{\textbf{权力流向(决策权归属)}} \\
    \cmidrule(lr){2-3}
    \textbf{权力形态(执行权分布)} & \textbf{决策权归中央} & \textbf{决策权归地方} \\
    \midrule
    \textbf{执行权分散}(行政分权) & 宋朝模式:枢密院、三衙、监军互相制衡 & 现代矩阵化条块治理;RPE 标尺竞争 \\
    \textbf{执行权集中}(行动集权) & 战时总动员;中央领导小组、指挥部 & 唐朝节度使模式 \\
    \bottomrule
  \end{tabularx}
\end{table}

四个格子可以逐一解读:

\begin{itemize}
  \item \textbf{决策集权 + 执行分散(宋朝模式)}。中央掌握最终决定权,但执行权被拆散到互不统属的机构,形成多重否决点(veto points)。\emph{是什么}:以安全换效率的保守设计。\emph{风险}:任何需要快速整合行动的任务都会瘫痪。
  \item \textbf{决策集权 + 执行集中(战时总动员)}。中央既决定目标,也把完成任务的完整工具箱交给统一的执行者。\emph{是什么}:动员效率最高的组合,常见于危机情境。\emph{风险}:高度依赖中央的信息与判断;若执行者脱离控制,则是最危险的组合。
  \item \textbf{决策分权 + 执行分散(矩阵化治理、标尺竞争)}。地方拥有执行自主权,但执行权被``条''``块''切割,或地方之间被置于互相竞争的标尺之下。\emph{是什么}:以竞争和制衡代替直接控制的现代技术。\emph{风险}:协调成本高,竞争产生负外部性。
  \item \textbf{决策分权 + 执行集中(唐朝模式)}。地方既决策又拥有完整工具箱。\emph{是什么}:任务效率最高的组合(免跨部门协调)。\emph{风险}:整合后的权力体量过大,一旦脱离中央控制即构成生存威胁。
\end{itemize}

\section{核心命题:两种真实需求的张力}

二维框架揭示了中国古代治理的核心矛盾——它从来不是简单的``中央要权''与``地方要权''之争,而是两个同样真实、同样合理的需求在组织设计上的冲突:

\begin{quote}
  \textbf{``中央为维护自身安全而拆散地方执行权''对``地方为完成治理任务而需要整合执行权''。}
\end{quote}

\begin{itemize}
  \item \textbf{系统外部压力}(战争、边患、自然灾害)要求力量集中、反应迅速,推动\textbf{行动集权};
  \item \textbf{系统内部稳定}(防止地方尾大不掉、威胁中央)要求权力制衡,推动\textbf{行政分权}。
\end{itemize}

古代中国没有能力同时满足这两个需求,只能在``一放就乱,一收就死''之间周期摇摆:外患加剧则打包权力,打包失控则拆散权力,拆散过度则国势不振,于是再次打包……循环往复两千年。第3章将展示唐、宋、明在这个循环中各自的位置。

% ==============================================================================
\chapter{以二维划分分析中国古代(唐宋明)的央地关系}

\section{引子:同一个困境的三个答案}

唐、宋、明三朝面对的是同一个困境,给出了三个不同的答案,并分别落在了二维矩阵的不同位置(表~\ref{tab:three}):唐朝选择打包(行动集权),宋朝选择拆散(行政分权),明朝试图在两者之间做有约束的临时切换。理解三朝,就是理解``拆''与``包''谱系上的三个位置及其代价。

\begin{table}[htbp]
  \centering
  \small
  \caption{唐、宋、明的二维定位与代价}
  \label{tab:three}
  \begin{tabularx}{\linewidth}{@{}lXX@{}}
    \toprule
    \textbf{王朝} & \textbf{方案} & \textbf{二维定位} \\
    \midrule
    唐朝 & 节度使:军、财、政打包给一人 & 决策分权 + 执行集中(行动集权) \\
    宋朝 & 三衙/枢密院/监军互相制衡 & 决策集权 + 执行分散(行政分权) \\
    明朝 & 督抚:中央特派的临时统筹者 & 中央决策 + 有约束的临时打包(时间切换) \\
    \bottomrule
  \end{tabularx}
\end{table}

\section{唐朝(618--907):行动集权的失控}

\textbf{是什么。} 唐朝在边境设置节度使,将军事指挥权、地方财政权、行政人事权``打包''赋予一人。

\textbf{为什么这么做。} 漫长的边境线面临突厥、回鹘、吐蕃等多方向军事压力。从组织理论看,节度使制是一种\textbf{任务驱动的权力整合}(task-driven power integration):决策者同时拥有资源和执行工具,不需要跨部门协调,因此效率极高——盛唐的开疆拓土正是其短期红利。

\textbf{卡在哪里。} 整合后的权力体量太大,一旦脱离中央控制就构成生存威胁。安史之乱(755--763)爆发时,中央能够调动的兵力远不如一个边镇节度使——不是因为中央穷,而是因为权力被打包给了地方。安史之乱后藩镇割据延续百余年,唐朝再未恢复对地方的完整控制。

\textbf{二维定位。} 决策分权 + 执行集中:象限 IV(行动集权)。\textbf{教训:打包而无制衡,必然尾大不掉。}

\section{宋朝(960--1279):行政分权的极致}

\textbf{是什么。} 赵匡胤以陈桥兵变上台,对武将拥兵自重有切肤之痛,于是做了四件事:杯酒释兵权,收回开国将领兵权;三衙(统兵)与枢密院(调兵)分离,互不统属;更戍法,军队定期换防,``兵不识将,将不识兵'';文官监军,以文御武。

\textbf{为什么这么做。} 内部安全优先:把执行权拆散,确保谁也不能威胁皇帝。这相当于在组织的每个关键节点设置多重否决点——任何人想做大事,都必须获得多个互不统属部门的同意。

\textbf{卡在哪里。} 内部安全的目标达到了:两宋319年无武将政变,文治达到高峰。但对外部威胁的响应速度被否决点严重拖慢:面对辽、金、蒙古的骑兵冲击,等枢密院的调令送到,战场形势已经变了。宋军屡战屡败,不是士兵不勇敢,而是组织设计不鼓励勇敢。最终以岁币换和平、以积弱换稳定,直至崖山。

\textbf{二维定位。} 决策集权 + 执行分散:象限 I(行政分权)。\textbf{教训:拆得过散,必然国势不振。}

\section{明朝(1368--1644):督抚制的折中}

\textbf{是什么。} 设立``督抚''——中央特派大员临时统筹数省军政。权力在需要时临时打包,但打包者是中央的人、打包有时间限制、打包后仍受多重监督(都察院、厂卫监视、任期限制)。

\textbf{为什么这么做。} 明朝认识到宋朝分权过度的低效,又不愿重蹈唐朝藩镇覆辙,试图在``拆''与``包''之间找第三条路。

\textbf{机制。} 督抚制是一种\textbf{有约束的任务型集权}(constrained task-oriented centralization):人事权(谁来做、做多久)牢牢控制在中央手中,同时给执行者一个完成特定任务所需的临时工具箱。效果上比宋能打、比唐稳定。

\textbf{卡在哪里。} 靠人事手段解决制度性问题有天然上限。督抚的临时性随任务常态化而褪色:频繁的边患与民变使督抚逐渐从临时差遣变成常设职位;一旦常态化,唐朝节度使的逻辑就开始悄悄复活。终明之世,地方军阀化苗头(如晚明左良玉)依然存在。

\textbf{二维定位。} 中央决策 + 有约束的临时打包:在象限之间做``时间切换''。\textbf{教训:人事手段可以调和张力,但不能根本解决;``临时''有滑向``永久''的内在趋势。}

\section{三朝小结:谱系而非类型}

\begin{table}[htbp]
  \centering
  \small
  \caption{三朝制度在``拆--包''谱系上的比较}
  \label{tab:compare3}
  \begin{tabularx}{\linewidth}{@{}lXXX@{}}
    \toprule
    & \textbf{唐:节度使} & \textbf{宋:制衡体系} & \textbf{明:督抚制} \\
    \midrule
    权力形态 & 打包(军财政一体) & 拆散(多重否决点) & 有约束的临时打包 \\
    制衡方式 & 几乎没有 & 过度制衡 & 人事 + 监察 + 任期 \\
    短期效果 & 边防高效,开疆拓土 & 内部安稳,文治高峰 & 比宋能打,比唐稳定 \\
    长期代价 & 安史之乱,藩镇割据 & 对外积弱,屡战屡败 & 督抚常态化,军阀化苗头 \\
    主要教训 & 行动集权无制衡 $\rightarrow$ 尾大不掉 & 行政分权过度 $\rightarrow$ 国势不振 & 人事折中治标不治本 \\
    \bottomrule
  \end{tabularx}
\end{table}

三朝的故事应当被读作\textbf{谱系而不是类型学}:唐、宋、明不是三种互斥的``制度类型'',而是同一张力在``拆--包''谱系上的三个位置。每个方案都有代价,历史没有标准答案。

\begin{center}
  \small
  \textbf{唐}(打包,失控) $\cdots$ \textbf{明}(有约束的临时打包) $\cdots$ \textbf{宋}(拆散,瘫痪)
\end{center}

更重要的是,古代中国缺少几个关键工具,才使这个循环无解:可靠的\textbf{信息}手段(``山高皇帝远'',中央无法验证地方上报的真伪)、\textbf{法治化的冲突裁决机制}(一切最终诉诸皇帝个人裁决)、独立的\textbf{审计与监察基础设施}。这三样工具的缺席,决定了古人只能在``拆''与``包''之间做零和切换。美国(第4章)与现代中国(第5章)的差异,正可以从这三个工具的有无来理解。

% ==============================================================================
\chapter{以二维划分分析美国联邦与州的央地关系}

\section{演化路径:一个持续两百多年的动态平衡}

美国的央地关系是演化出来的,不是设计出来的(表~\ref{tab:us-path}):

\begin{table}[htbp]
  \centering
  \small
  \caption{美国联邦--州关系的演化路径}
  \label{tab:us-path}
  \begin{tabularx}{\linewidth}{@{}lXX@{}}
    \toprule
    \textbf{阶段} & \textbf{权力格局} & \textbf{核心事件} \\
    \midrule
    邦联时期(1776--1787) & 州权至上,联邦形同虚设 & 《邦联条例》下联邦无权征税、无力偿还债务 \\
    宪法确立(1787) & 联邦权力列举,剩余归州 & 康涅狄格妥协;第十修正案;``精心设计的模糊'' \\
    南北战争(1861--1865) & 以战争回答州能否退出 & 约62万人阵亡,确立联邦不可分裂 \\
    二元联邦主义(19世纪) & 联邦与州各有独立管辖范围 & ``分层蛋糕'':井水不犯河水 \\
    合作联邦主义(1930s--) & 联邦拨款渗透州的政策领域 & 罗斯福新政;``大理石蛋糕'':边界模糊化 \\
    新联邦主义(1970s--) & 部分权力归还各州 & 钟摆回摆,寻求新平衡 \\
    \bottomrule
  \end{tabularx}
\end{table}

\section{制度特色:三种关键设计}

\begin{itemize}
  \item \textbf{法治化裁决}。联邦最高法院是央地纠纷的最终权威。1819年``麦卡洛克诉马里兰''案确立联邦隐含权力(implied powers),是央地边界由司法裁决而非个人意志决定的标志性案例。
  \item \textbf{复合共和(双重主权)}。联邦政府直接面对公民,拥有自己的执法系统(如 FBI),不依赖州政府执行联邦法律。这与中国传统模式中``地方是中央的派出和执行机构''形成结构性差异。
  \item \textbf{动态平衡}。通过时代演变、政治过程与司法审查,持续调整央地权力边界,而非追求一劳永逸的``最终答案''。
\end{itemize}

\section{用二维框架看美国的央地关系}

把二维框架套到美国制度上,可以看到几个关键观察(表~\ref{tab:us-matrix}):

\begin{table}[htbp]
  \centering
  \small
  \caption{美国制度环节的二维定位}
  \label{tab:us-matrix}
  \begin{tabularx}{\linewidth}{@{}lXX@{}}
    \toprule
    \textbf{制度环节} & \textbf{二维定位} & \textbf{说明} \\
    \midrule
    州在其保留领域 & 决策分权 + 执行集中 & 州有自己的宪法、议会、法院、执法系统,拥有完整工具箱 \\
    联邦直接执法领域 & 决策集权 + 执行集中 & 联邦决定目标并直接执行(FBI、联邦法院),不依赖州 \\
    联邦拨款附带条件、州执行 & 决策集权 + 执行分散 & 联邦设定规则,执行分散在各州,类似``条''对``块''的渗透 \\
    州际竞争与``用脚投票'' & 决策分权 + 分散制衡 & 州有自主权,但被居民迁徙(Tiebout 排序)与横向比较约束 \\
    \bottomrule
  \end{tabularx}
\end{table}

由此可以得到四个一般性观察:

\textbf{观察一:``流向''被宪法固定为程序问题。} 美国模式的根本差异不在于``谁更集权'',而在于央地冲突的\textbf{裁决机制}:权力边界由宪法规定、由最高法院解释,而不是由单一权威拍板。``流向''的每一次调整都经过规则化的博弈,而非一次性摊牌。

\textbf{观察二:州层面是``有约束的行动集权''。} 州在自己的保留领域内拥有决策与执行的完整工具箱(象限 IV 的结构),但它处于三重约束之下:联邦宪法与最高法院、横向州际竞争、以及退出权被内战永久取消。也就是说,美国把唐朝节度使式的``打包''放进了法治化的笼子里——约束来自\textbf{制度}而非\textbf{人事}。

\textbf{观察三:州际竞争即标尺竞争。} 中央不直接控制各州做什么,但通过相对比较间接约束:居民和企业可以``用脚投票''(Tiebout 1956),各州绩效互相成为标尺。这与当代中国的晋升锦标赛在机制上同源(见 5.1.4 节)——都是``给地方完整执行权,但让多个地方互相竞争''。

\textbf{观察四:矛盾从未消失,只是被容纳。} 2024 年初得州与联邦政府在移民问题上的武装对峙说明,央地矛盾不可能被``一劳永逸地解决''——它是一个需要被持续管理的动态张力。美国用两个工具管理它:宪法作为博弈规则,最高法院作为规则解释者。这不完美,但提供了一个比``皇帝拍板''更稳定、更可预期的冲突解决机制。

\section{与中国传统模式的对比}

\begin{table}[htbp]
  \centering
  \small
  \caption{央地冲突裁决的结构性差异}
  \label{tab:us-cn}
  \begin{tabularx}{\linewidth}{@{}lXX@{}}
    \toprule
    \textbf{维度} & \textbf{美国模式} & \textbf{中国传统模式} \\
    \midrule
    权力来源 & 联邦权力来自宪法授权,剩余归州 & 所有权力最终来源于中央授权 \\
    冲突裁决 & 联邦最高法院(司法裁决) & 皇帝/中央(政治裁决) \\
    地方独立性 & 州有宪法、议会、法院、执法系统 & 地方是中央的派出和执行机构 \\
    联邦/中央执法 & 联邦机构直接行动,不依赖州 & 依赖地方执行中央政令 \\
    \bottomrule
  \end{tabularx}
\end{table}

需要强调的是,表~\ref{tab:us-cn} 是制度机制层面的比较,不是优劣评判:美国方案以决策链条长、协调成本高为代价,换取规则化的博弈框架;这同样是 trade-off,不是免费的答案。本报告不做``应向美国学习''的规范性建议。

% ==============================================================================
\chapter{以二维划分分析中国当代的央地关系}

\section{现代管理手段:从零和切换到同时存在}

古代中国在``拆''与``包''之间做的是零和切换:要安全就拆散,要效率就打包,二者不可兼得。当代中国的制度创新之处在于,凭借古代没有的三个工具——矩阵化的组织设计、穿透式的信息技术、任务型的组织形态——让``拆''与``包''\textbf{不再是非此即彼的时间切换,而是可以在不同维度上同时存在};此外还有第四个手段——相对绩效评估(RPE),以竞争代替直接控制。这四个手段构成了当代央地管理的工具箱。本节先逐一展开四个手段(5.1.1--5.1.4),随后两节用南方电网与土地财政两个案例检验它们的组合运作,最后一节小结当代中国在二维矩阵中的位置。

\subsection{条块关系与矩阵化管理}

\textbf{是什么。} 当代中国的治理结构通常被描述为``条块结合'':``条''(tiao)是中央部委到地方对口部门的垂直职能系统(公安部 $\rightarrow$ 省公安厅 $\rightarrow$ 市公安局);``块''(kuai)是以行政区划为单位的地方政府对本级各部门的横向统合。``条块结合''即双重领导:省公安厅同时受公安部和省政府领导。

\textbf{为什么重要。} 矩阵化管理的本质,是在\textbf{不同维度}分配不同类型的权力:公安厅长的人事和经费由公安部垂直管理(条),但他的日常行动受省政府的横向协调(块)。制度层面上权力被拆散,执行层面上可以临时打包——两者同时存在。这比宋朝的多头制衡更精细:宋朝是在\textbf{同一个维度}(军事指挥)上制造否决点,矩阵化则是在\textbf{不同维度}上分配不同权力。

\textbf{卡在哪里。} 双重领导天然产生模糊地带:条块责任边界不清时,容易出现互相推诿或政出多门;条块之间的张力随政策周期此消彼长。

\subsection{穿透式监管:信息技术改变信息不对称}

\textbf{是什么。} 古代治理最大的信息瓶颈是``山高皇帝远'':中央无法验证地方上报信息的真实性,信息是稀缺的、向上的、容易被扭曲的。现代信息技术——实时数据采集、遥感监测、卫星监控、财务电子化、在线审计——让中央可以穿透层级直接看到基层,部分打破了地方的``信息黑箱''。

\textbf{为什么重要。} 如果中央能在不收回执行权的情况下保持对地方行为的感知能力,那么``打包''权力的风险就大幅降低:前线指挥官的集权,是在透明笼子里的集权,随时可以被中央察觉和纠正。这正是南方电网等``区域行动集权''安排得以成立的前提之一(见 5.2 节)。

\textbf{卡在哪里。} 这一机制在包内的证据状态被标为\textbf{未定}(unresolved):(a)信息获取能力 $\neq$ 信息处理能力——中央可能被数据淹没;(b)地方会发展新的规避策略(数据造假升级);(c)监督的``最优密度''可能是倒U型而非单调递增——信息透明也可能诱使中央过度干预,扼杀地方能动性。这些问题尚无系统的经验研究(见 6.3 节)。

\subsection{任务型组织:临时集权的制度化}

\textbf{是什么。} ``领导小组''``指挥部''等因事而聚、事毕则散的临时性机构:面临重大危机(抗疫、抢险、重大攻关)时,中央迅速组建指挥部,将相关权力``打包''给前线负责人;任务完成后机构解散,权力回流。

\textbf{为什么重要。} 这试图实现唐朝节度使的执行效率,同时避免其永久化。前提是现代国家有足够的制度基础设施(公文系统、人事档案、财政审计)来确保``临时''真的可以是临时的。

\textbf{卡在哪里。} 历史上``临时''滑向``永久''的案例比比皆是——明朝督抚就是最好的先例,当代也不乏存续数十年的``领导小组''。``因事而聚,事毕则散''的成立条件尚未被清晰说明(见 6.2 节)。\src{IP}

\subsection{相对绩效评估(RPE):竞争的间接控制}

\textbf{是什么。} 相对绩效评估(relative performance evaluation)是委托--代理理论的核心工具,由 Holmstr\"om(1982)形式化。基本逻辑:评估代理人时不要只看绝对产出,要拿他和同行比——所有代理人都承受共同的随机冲击(市场、政策、气候),减去同行平均表现,剩下的才是个人努力的信息。形式化:

\begin{equation}
  \text{评估信号} \;=\; y_i - \rho \cdot \bar{y}_{\text{peer}}
\end{equation}

其中 $y_i$ 是代理人 $i$ 的产出,$\bar{y}_{\text{peer}}$ 是同行平均产出,$\rho$ 是共同冲击的相关程度($\rho$ 越高越应依赖相对评价)。

\textbf{为什么重要:RPE 是行政分权的一种实现工具。} 在二维框架中,宋朝的答案是``拆散执行权''——代价是能力瘫痪。RPE 提供了替代路径:\textbf{给地方完整执行权,但让多个地方互相竞争}。地方有自主权(不会``死''),竞争压力约束行为(不会``乱'')。中国地方官员的晋升锦标赛(周黎安 2007,\src{IP} 争议状态)正是 RPE 的典型应用:中央通过比较各省相对 GDP 增速决定升迁,比直接考核``是否努力工作''更省监督成本。

\textbf{卡在哪里:竞争制造新的央地矛盾。} (a)负外部性:地方拒绝跨域合作(跨区域污染治理、基础设施衔接);(b)标尺操纵:地方优化统计而非优化实际;(c)标尺同质化:所有地方追逐同一指标,丧失地方多样性(``千城一面'')。此外,穿透式监管(5.1.2 节)与 RPE 存在交互:能直接观测时,为何还要通过同行比较间接推断?但信息技术也让 RPE 更精细(多维指标的相对位置)——这是一个尚未回答的经验问题。

\section{案例一:南方电网——区域行动集权 + 中央决策集权}\label{sec:csg}

\subsection{是什么:一个功能需求驱动的新制度实体}

2002 年 2 月,国务院印发《电力体制改革方案》(国发〔2002〕5 号),确立``厂网分开、竞价上网''。原国家电力公司被拆分:发电资产重组为五大发电集团;电网资产一分为二——国家电网覆盖 26 个省区,\textbf{中国南方电网有限责任公司}覆盖广东、广西、云南、贵州、海南五省区,2002 年 12 月 29 日挂牌。南方电网是唯一总部不设在北京的电网央企(总部在广州),也是唯一股权结构中包含地方国有资本的电网央企:\src{SC}

\begin{table}[htbp]
  \centering
  \small
  \caption{南方电网股权结构(2002 年组建时)}
  \label{tab:csg-equity}
  \begin{tabularx}{\linewidth}{@{}lXl@{}}
    \toprule
    \textbf{股东} & \textbf{性质} & \textbf{持股比例} \\
    \midrule
    国务院国资委 & 中央 & 51\% \\
    广东恒健投资控股有限公司 & 广东省属 & 25.57\% \\
    中国人寿保险(集团)公司 & 中央金融企业 & 21.3\% \\
    海南省发展控股有限公司 & 海南省属 & 2.13\% \\
    \bottomrule
  \end{tabularx}
\end{table}

南网的雏形不是行政划分的产物,而是\textbf{工程任务驱动的产物}:2000 年北戴河会议讨论``十五''规划时,广东提出新增 1000 万千瓦电力装机的诉求;中央领导人反对在广东境内集中新建火电,主张用西电东送解决缺口;最终达成``拼盘方案''——云贵桂到广东的西电东送南线通道 + 三峡到广东的``三广线''直流工程 + 湖南鲤鱼江电厂扩建,合计正好 1000 万千瓦。涉及的五省区电网互联,后来就构成了南方电网的地理基础。\footnote{北戴河会议讨论细节来自包内二手转述,\src{NV},引用前需交叉验证。}

\subsection{二维定位:象限 II 的区域变体}

南方电网的成立属于\textbf{``区域行动集权 + 中央决策集权''的组合}:

\begin{itemize}
  \item \textbf{决策权归中央}:中组部任命领导班子;国资委行使控股权(51\%);重大投资与电价受国家发改委、国家能源局约束;
  \item \textbf{执行权打包}:五省区电网统一调度,西电东送由南网统一规划建设运营,跨省交易在公司内部协调。
\end{itemize}

换言之,南网在南方五省区内实现了类似唐朝节度使的\textbf{行动效率}(统一的执行工具箱),同时通过现代制度基础设施保持\textbf{中央控制}。与唐、宋的对比见表~\ref{tab:csg-compare}:

\begin{table}[htbp]
  \centering
  \small
  \caption{唐节度使、宋制衡体系与南方电网的比较}
  \label{tab:csg-compare}
  \begin{tabularx}{\linewidth}{@{}lXXX@{}}
    \toprule
    & \textbf{唐:节度使} & \textbf{宋:制衡体系} & \textbf{南方电网} \\
    \midrule
    打包程度 & 极高(军财政一体) & 极低(互相拆散) & 中等偏高(五省电网运营权打包) \\
    中央控制方式 & 人事任命,可世袭拥兵 & 多重否决点 & 股权 + 人事 + 行业监管 \\
    绩效 & 短期高效,长期失控 & 长期稳定,对外低效 & 20+ 年运营,效率较高且未失控 \\
    关键差异 & 无外部制衡 & 过度制衡致能力瘫痪 & 制度化多层制衡(股权、审计、监管、信息公开、对标竞争) \\
    \bottomrule
  \end{tabularx}
\end{table}

南网之所以比唐朝节度使更可持续,不是因为中央更``信任''南网,而是因为现代国家有唐朝不具备的制度基础设施:股权控制(随时可换人)、财务审计(信息不对称大幅降低)、行业监管(偏差可及时检测)、以及与国网并存形成的\textbf{对标比较}(yardstick competition)环境——中央可以通过比较两家电网的相对绩效,获得比单一垄断者时代更准确的信息。

\subsection{中央的设计逻辑(含待核验的隐性期待)}

从公开政策与制度安排可梳理出四条逻辑:\src{SC}

\begin{enumerate}
  \item \textbf{跨省协调的执行主体}:西电东送跨越多个省级行政区,需要超越各省利益的区域统筹者,把省际谈判内部化为公司内部调度;
  \item \textbf{可比竞争的试验场}:双头结构为中央提供相对绩效评估的信息基础(与 5.1.4 节的 RPE 逻辑同源);
  \item \textbf{广东因素的纳入式妥协}:广东电网原属地方资产,简单收归中央会遭遇阻力;让广东省属资本持股 25.57\%、总部设在广州,是把地方利益\textbf{制度化嵌入}央企治理,以股权嵌入实现利益绑定;
  \item \textbf{为后续改革保留制度弹性}:保留两个可比主体,使输配分离、调度独立等后续改革更易操作。
\end{enumerate}

第 2--4 条中含``中央隐性期待''成分,属于基于制度安排的反向推理(\src{IP}),并非基于一手政策文件或决策者回忆的直接证据。

\subsection{卡在哪里:三层次张力与模式局限}

南网运行 20 余年后,制度设计的优势已获验证,但新的张力正在浮现,可归纳为三个层次:\src{SC}(问题事实) + \src{IP}(严重程度判断)

\begin{enumerate}
  \item \textbf{功能张力(技术--经济层)}:新能源的间歇性对``西部发电 $\rightarrow$ 东部用电''的传统调度构成根本挑战(2023--2024 年云南干旱使全国水电出力下降 23\%,南网被迫依赖广东本地煤电);新型电力系统建设需要巨额资本支出,而输配电价``准许成本 + 合理收益''监管压制回报;
  \item \textbf{利益张力(央--地--省际层)}:西电东送的电价与电量安排中,送端省有``穷省补贴富省''的隐忧,受端省有``我出钱修通道''的立场——这是南网制度设计最核心的考验。若省际矛盾越来越需要中央出面裁断,``内部化''设计的有效性就值得重新审视;
  \item \textbf{治理张力(央企属性 vs 混合股权)}:广东 25.57\% 股权在关键决策中究竟是实质治理参与,还是``仪式性共治''?缺乏公开研究。这直接决定``纳入广东利益''的制度设计是否真实存在。\src{NV}
\end{enumerate}

模式的局限同样明确:可推广性存疑(高度依赖广东的经济体量);双头竞争的信息揭示功能可能被高估(两家公司不在同一地理市场竞争);制度持续性未经重大冲击检验——唐朝节度使在初期也运行良好。

\section{案例二:土地财政——行政分权的财政版}\label{sec:land}

\subsection{是什么}

``土地财政''是对中国地方政府高度依赖土地相关收入维持运转的现象的描述,包括三部分:土地出让金(地方政府出让土地使用权的一次性收入,通常占地方本级财政收入的 40\%--60\%)、土地相关税收、以及土地抵押融资(城投公司以土地为抵押借款,不直接入账但构成实际支出能力——也是最危险的部分)。巅峰时期(2020--2021 年)全国土地出让金超过 8 万亿元;部分三四线城市的土地出让金占地方综合财力比重超过 100\%。

\subsection{形成:1994 年分税制之后的制度性求生}

1990 年代初中央财政陷入危机:中央财政收入占全国财政收入的比重从 1980 年代的约 40\% 降至 1993 年的 22\%,全国财政收入占 GDP 比重从 1978 年的 31\% 降至 12\%。1993 年国务院副总理朱镕基不得不亲赴各省``化缘''。1994 年分税制改革将最丰厚的税种(增值税的 75\%、消费税的全部等)划归中央,中央财政收入比重迅速回升至 55\% 左右——这是改革开放后最重大的一次中央集权。

但改革留下了结构性后遗症:财权上收,事权未相应上收(教育、医疗、基建支出仍压在地方),转移支付体系当时尚未建立。用一句话概括:\textbf{``钱我拿走了,但活儿还是你干。''}地方不能印钞、不能举债(《预算法》当时禁止)、不能设税,却发现了制度裂缝:\textbf{土地出让金被规定为地方固定收入,不参与中央分成}。1998 年住房商品化改革准备了需求端,2002 年经营性用地``招拍挂''制度点燃了地价——土地财政由此爆发。

\subsection{二维定位:拆散财源与打包自救}

\begin{table}[htbp]
  \centering
  \small
  \caption{土地财政各阶段的二维定位}
  \label{tab:land-matrix}
  \begin{tabularx}{\linewidth}{@{}p{2.6cm}XX@{}}
    \toprule
    \textbf{阶段} & \textbf{中央行为} & \textbf{地方行为与二维定位} \\
    \midrule
    1994 分税制 & 财权上收(决策集权) & 事权保留;地方财源被拆散——中央的\textbf{行政分权} \\
    2003--2016 膨胀 & 默许 & 土地出让 + 抵押 + 基建打包使用——地方的\textbf{行动集权}(制度裂缝中的自救) \\
    2016--2021 收紧 & ``房住不炒''、三条红线、集中供地 & 地方土地融资工具被拆散——中央再\textbf{行政分权} \\
    2022-- 收缩 & 消费税下划、扩大转移支付 & 地方``戒断反应''——尚未形成新均衡 \\
    \bottomrule
  \end{tabularx}
\end{table}

用二维框架看,土地财政完整演绎了``一收就死,一放就乱''的财政版:分税制拆散地方财源(行政分权) $\rightarrow$ 地方财政缺口 $\rightarrow$ 地方在制度裂缝中打包土地工具自救(行动集权) $\rightarrow$ 中央默许(因为土地收入推动城市化和 GDP 增长、减轻转移支付压力) $\rightarrow$ 杠杆膨胀、房价畸高、金融风险 $\rightarrow$ 中央不得不收(三条红线等再行政分权) $\rightarrow$ 土地市场崩溃、地方财政承压 $\rightarrow$ 中央被迫再放(消费税下划、兜底地方债)——新一轮循环尚未走完。

\textbf{与宋朝的类比}是深刻而具体的:宋朝中央为自身安全拆散了地方的军事执行权,前线无力作战,被迫走偏门(岁币求和);分税制后中央为集中财力拆散了地方的财源,地方无力正常运转,被迫走偏门(土地财政)。但现代地方的自救机制比宋朝危险得多:宋朝地方无能打仗至少不会造成系统性金融风险,而城投债违约会经金融系统传导全国。

\subsection{代价与现状}

代价是结构性的:地方从``服务市民''转向``经营城市''(鬼城、摊大饼式扩张);拆迁冲突成为 2000--2010 年代最激烈的社会矛盾之一;代际财富分配被房价扭曲;IMF 估计 2023 年地方政府隐性债务(主要是城投债)约 50--60 万亿元,相当于 GDP 的 40\%--50\%。

2021 年恒大违约引发房企连锁爆雷,2022 年人口首次负增长,土地出让金从 2021 年的约 8.7 万亿元断崖式下跌(2023 年约 5.8 万亿元,2024 年估算约 4 万亿元)。地方财政出现``戒断反应'':部分地区公务员降薪、三四线城市公交停运、罚没收入激增。替代方案(消费税下划、房产税、扩大转移支付、地方债显性化与中央兜底)各有难处,尚未形成新均衡。

\emph{证据提示}:``中央默许土地财政''的判断属于制度推理(\src{IP}),来自对制度安排的反向推理,未对照决策层内部讨论的一手文献;土地出让金数据以官方公布为准,城投债隐性债务规模引用 IMF 估算。\src{NV}

\section{小结:当代中国在二维矩阵中的位置}

当代中国在二维矩阵中的位置\textbf{不是一个点,而是一组动态组合}:

\begin{itemize}
  \item 常规条块治理:决策集权 + 条块分散执行(象限 I 与 III 的混合);
  \item 领导小组、指挥部:象限 II 的\textbf{临时}进入(因事打包,事毕则散);
  \item 南方电网:象限 II 的\textbf{区域常设}变体(区域行动集权 + 中央决策集权);
  \item 晋升锦标赛、RPE:决策分权 + 分散竞争(象限 III 的标尺约束);
  \item 土地财政时代的地方行为:财政维度上的象限 IV(地方打包自救),已被新一轮行政分权部分拆散。
\end{itemize}

历史循环被\textbf{部分}打破的条件,正是唐、宋、明所没有的三个工具:矩阵化组织(在不同维度分配权力)、信息技术(降低监督成本)、任务型组织(制度化临时集权)。但新工具带来了新张力:监督与能动性的权衡、竞争与合作的冲突、``临时''与``永久''的漂移——它们构成了第6章的开放问题。

% ==============================================================================
\chapter{开放性问题}

本章汇总本报告框架尚未回答、且据包内判断在学术上也尚未被充分回答的问题。编号沿用包内缺口表(G1--G5),并补充案例研究中浮现的新问题。

\section{G1:二维框架本身——形式化与文献对齐}

现有文献中,``集权/分权''通常被当作一维连续变量处理。本报告主张宋朝模式(决策集权 + 执行分散)与唐朝模式(决策分散 + 执行集中)在一维轴上可能都落在``中间'',但结构完全不同。\textbf{待验证}:政治学与公共行政文献中关于 decentralization 的多维度量方案(如政治/行政/财政分权的标准区分)是否已经覆盖了``执行权拆散 vs 打包''这一维度?

\begin{itemize}
  \item 若缺口成立:可发展为概念性论文或测量方法论文——大量用一维度量分权并回归治理绩效的跨国研究可能遗漏了关键的结构维度;
  \item 若缺口不成立:应将本包语言对齐到已有学术术语,保留其分析启发价值。
\end{itemize}

\section{G2:临时集权永久化的机制条件}

``领导小组''``指挥部''声称``因事而聚,事毕则散'',但历史上大量``临时''制度都永久化了:唐朝节度使、明朝督抚都是先例。\textbf{开放问题}:在什么条件下临时集权会滑向永久?什么制度设计可以防止这种蜕变?现有文献以案例研究为主,缺乏统一分析框架。最小验证方式:系统收集不少于 10 个``临时机构 $\rightarrow$ 永久化''与``临时机构 $\rightarrow$ 确实解散''的对比案例,用过程追踪提取差异条件。制度变迁理论中的路径依赖(Pierson 2004)与制度叠加(Thelen 2004)可提供理论资源。

\section{G3:信息技术对央地关系的因果效应}

大量讨论声称大数据、实时监控、在线审计正在改变央地关系,但缺乏严格经验研究回答:(a)信息技术是否真的降低了中央的监督成本?(b)地方是否发展出了新的规避策略?(c)对治理绩效的净效应是正还是负?若信息技术确实改变了委托--代理方程的基本参数,大量基于传统信息不对称假设的理论需要修正;若效果被夸大,政策资源可能被错误配置。最小验证方式:寻找信息技术在央地关系中的``自然实验''(如某在线审计系统在不同省份/领域的渐次推广),做双重差分(DID)分析。

\section{G4:最优分权度的识别}

财政联邦主义理论(Oates 1972 的分权定理)给出了分权最优程度的原则(分权收益与外溢效应、规模经济相权衡),但缺乏可操作的经验识别方法:分权程度本身是内生的(国家选择分权程度,而非随机分配),跨国面板回归面临严重内生性。若不解决识别问题,``最优分权度''就始终是理论概念而非政策工具。这是一个方法困难多于理论困难的问题。

\section{G5:跨学科映射的分析价值}

将央地关系映射为层级控制系统(监督 = 反馈回路,地方自由裁量权 = 增益参数,信息扭曲 = 传感器噪声)在概念上有吸引力,但\textbf{类比本身不是证据}。需要证明该映射能产生超出政治学现有方法的洞察。最小验证方式:构建极简模型(控制论或 agent-based),考察``一放就乱、一收就死''的周期性行为能否在两个学科框架下得到实质不同的解释。风险是低维模型只是在``重新描述''已知现象。

\section{案例研究带来的新问题}

\begin{itemize}
  \item \textbf{南网模式的可推广性}:区域行动集权 + 中央决策集权在欠发达区域是否成立?广东 25.57\% 股权的实质治理影响是什么(治理参与还是仪式性共治)?
  \item \textbf{土地财政之后的地方财源}:消费税下划与房产税在政治上各有多大可行性?``数据财政''``碳财政''等变体能否替代土地财政?人口负增长意味着所有依赖``人''的地方税制都有极限——这是分税制修补解决不了的问题。央地财政关系是否存在根本重构(而不仅是修补)的可能?
  \item \textbf{非正式制度的角色}:央地关系中大量非正式的协商、谈判与变通,是正式制度的``润滑剂''还是``腐蚀剂''?非正式制度天然难以观测,这是方法瓶颈。
  \item \textbf{跨国比较的实证基础}:目前多数央地关系研究是单一国家案例,缺乏可比的跨国量化指标。
\end{itemize}

\section{缺口汇总}

\begin{table}[htbp]
  \centering
  \small
  \caption{开放问题/研究缺口汇总}
  \label{tab:gaps}
  \begin{tabularx}{\linewidth}{@{}lXll@{}}
    \toprule
    \textbf{编号} & \textbf{问题} & \textbf{类别} & \textbf{方法难度} \\
    \midrule
    G1 & 二维框架的形式化与文献对齐 & 疑似缺口(suspected) & 低 \\
    G2 & 临时集权永久化的条件 & 确认缺口(confirmed) & 高 \\
    G3 & 信息技术的因果效应 & 确认缺口(confirmed) & 高 \\
    G4 & 最优分权度的识别 & 疑似缺口(suspected) & 很高 \\
    G5 & 控制论/系统科学映射的价值 & 弱缺口(weak) & 中 \\
    \bottomrule
  \end{tabularx}
\end{table}

\emph{延伸方向}:包内已将上述缺口转化为四个候选项目(PC1 控制论模型、PC2 历史案例结构化比较、PC3 二维框架概念论文、PC4 agent-based 模型,均为 provisional 状态,详见 \texttt{22\_project\_candidates.md}),本报告不再展开。

% ==============================================================================
\chapter*{结语}
\addcontentsline{toc}{chapter}{结语}

把整条线索收拢,可以得到四个层次的结论:

\begin{enumerate}
  \item \textbf{术语层面}:不要混淆``权力流向''与``权力形态''。权力流向中央不等于集权(可以是拆散),权力流向地方不等于分权(可以是打包)。唐、宋的教训分别对应这两个方向上的误读。
  \item \textbf{机制层面}:治理的永恒张力是——系统外部压力要求行动集权,系统内部稳定要求行政分权。两个需求都真实、都合理,但在组织设计上互相冲突。古代中国只能在``一放就乱,一收就死''之间循环。
  \item \textbf{比较层面}:历史没有标准答案。唐太松、宋太紧、明在中间摇摆——每个方案都有代价。美国给出了另一种解法:不是靠人,而是靠制度与司法,把不可消除的张力转化为可管理的规则化博弈。这是启示,不是可以直接照搬的答案。
  \item \textbf{工具层面}:现代工具部分改变了游戏。矩阵化组织 + 信息技术 + 任务型组织,让``既能集、又能散''的非零和博弈在当代中国初步成为可能;南方电网展示了``区域行动集权 + 中央决策集权''的常设形态,土地财政则警示了财政维度上``拆''与``包''循环的现代重演。最高明的治理,不是纠结于权力该``流回中央''还是``流向地方'',而是搞清楚:在制度设计上让权力分散到足以制衡腐败和分裂,在具体执行上让权力集中到足以应对危机和挑战。这条路是否可持续、新工具的新张力如何演化,正是第6章开放问题的意义所在。
\end{enumerate}

% ==============================================================================
\appendix

\chapter{材料来源}

本报告整合自以下包内文件(路径相对 \texttt{political\_science/central\_local\_relations/}):

\begin{itemize}
  \item \texttt{conflict.md}——种子笔记:``央地矛盾、集权与分权的治理逻辑'';
  \item \texttt{00\_beginner\_guide.md}——人话入门指南;
  \item \texttt{01\_research\_overview.md}——领域总述报告;
  \item \texttt{02\_frontier\_compass.md}——边界罗盘(内圈/中圈/前沿/伪缺口);
  \item \texttt{11\_term\_disambiguation.md}——术语辨析与概念边界契约;
  \item \texttt{12\_field\_map.md}——问题链与方法谱系;
  \item \texttt{17\_claims\_evidence\_ledger.md}——论断--证据台账(C1--C8);
  \item \texttt{20\_gap\_table.md}——研究缺口表(G1--G5);
  \item \texttt{22\_project\_candidates.md}——候选研究项目(PC1--PC4);
  \item \texttt{90\_QA.md}——问答缓冲区(含 RPE、法团主义同构、土地财政等讨论);
  \item \texttt{case\_studies/southern\_power\_grid.md}——南方电网案例研究;
  \item \texttt{case\_studies/land\_finance.md}——土地财政案例研究。
\end{itemize}

\chapter{证据分层图例与关键标注点}

\section{图例}

\begin{itemize}
  \item \src{SC} source\_checked:有公开可查证来源支持的基本事实(如南网的股权结构、成立时间、土地出让金官方数据);
  \item \src{WD} working\_definition:包内工作定义或分析性构造,尚未与学术文献完全对齐(如``行政分权''``行动集权''的用法);
  \item \src{IP} interpretation:基于制度安排的反向推理,非一手证据(如中央的``隐性期待''、``中央默许土地财政'');
  \item \src{NV} NEEDS\_VERIFICATION:在当前条件下无法核验、引用前需重新查证的条目。
\end{itemize}

\section{关键标注点汇总}

\begin{itemize}
  \item 2.4 节:``行政分权/行动集权''为工作定义,与 deconcentration 标准含义的差异;
  \item 5.1.3 节:``任务型组织能避免永久化''属乐观主张,历史反例(明督抚)存在;
  \item 5.1.4 节:晋升锦标赛(周黎安 2007)在学术界存在争议(contested);
  \item 5.2 节:中央拆分决策中的``隐性期待''部分、2000 年北戴河会议细节、广东恒健在南网董事会的实质表决权,均属 \src{IP}/\src{NV};
  \item 5.3 节:``中央默许土地财政''属 \src{IP};城投债隐性债务规模引用 IMF 估算;
  \item 5.1.2/6.3 节:信息技术改变央地关系的论断证据状态为``未定''(unresolved)。
\end{itemize}

\chapter{参考文献(包内记录,部分待核验)}

\emph{说明}:以下条目均来自包内已有记录,本报告未新增引文。包内未给出完整题名的条目只列作者与年份,标注``题名待核验'';引用前请按作者--年份检索原文。政策文件与行业报告据包内记录列出。

\begin{itemize}
  \item Oates, W. E. 1972. \emph{Fiscal Federalism}.(财政联邦主义,分权定理)
  \item Tiebout, C. M. 1956. ``用脚投票''模型。(题名待核验)
  \item Kornai, J. 1986. 软预算约束。(题名待核验)
  \item Holmstr\"om, B. 1982. 相对绩效评估的形式化。(题名待核验)
  \item Hayek, F. A. 1945. 知识的分散性与决策权配置。(题名待核验)
  \item Seabright, P. 1996. 分权与问责的不完全契约模型。(题名待核验)
  \item Brennan, G., \& Buchanan, J. M. 1980. 竞争性政府的财政约束。(题名待核验)
  \item Bueno de Mesquita, B. et al. 2003. 选举人团理论(Selectorate Theory)。(题名待核验)
  \item Weingast, B. R.;Qian, Y., \& Roland, G. 1990 年代. 市场保护型联邦主义。
  \item Schmitter, P. C. 1974. 法团主义的经典定义。(题名待核验)
  \item Pierson, P. 2004. 制度变迁中的路径依赖。(题名待核验)
  \item Thelen, K. 2004. 制度变迁的 layering 理论。(题名待核验)
  \item 周黎安. 2007. 晋升锦标赛。(题名与出处待核验;状态:contested)
  \item 荣敬本等. 1990 年代. 压力型体制。
  \item 渠敬东等. 2010 年代. 项目制。
  \item 郑永年. 行为联邦制(Behavioral Federalism)。
  \item Lewis, J. I. 2023. \emph{Cooperating for the Climate}. MIT Press.
  \item Kahrl, F., Lin, J., Liu, X., \& Hu, J. 2021. ``Sunsetting coal power in China.'' \emph{iScience}, 24(9).
  \item Curtis, S., \& Klaus, I. 2024. \emph{The Belt and Road City}. Yale University Press.
  \item Zhang, A. H. 2024. \emph{High Wire: How China Regulates Big Tech and Governs Its Economy}.
  \item Pollitt, M., Yang, H., \& Chen, H. 2018. 中国电力体制改革的国际比较。\emph{Cambridge Working Papers in Economics}.
  \item 国务院. 2002. 《电力体制改革方案》(国发〔2002〕5 号).
  \item 中共中央、国务院. 2015. 《关于进一步深化电力体制改革的若干意见》(中发〔2015〕9 号).
  \item 国家发改委. 2020. 《省级电网输配电价定价办法》(修订版).
  \item 国家发改委、国家能源局. 2022. 《关于加快建设全国统一电力市场体系的指导意见》.
  \item 国家能源局. 2023. 《新型电力系统发展蓝皮书》.
  \item IEA. 2023. \emph{Electricity Market Report 2023}.
  \item \emph{The Economist}. 2023-11-27. ``Will China save the planet or destroy it?''
  \item 中国电力企业联合会(CEC). 各年度电力供需分析报告.
  \item IMF. 2023. 中国地方政府隐性债务估算(约 50--60 万亿元).
\end{itemize}

\end{document}
